\documentclass[a4paper,11pt]{article}
\usepackage[a4paper,margin=2cm]{geometry}
\usepackage[T1]{fontenc}
\usepackage[utf8]{inputenc}
\usepackage[ngerman,english]{babel}
\usepackage{lmodern}
\usepackage[hidelinks]{hyperref}
\title{Freier Akkusativ: adverbialer Akkusativ der Zeit und des Maßes}\date{}\begin{document}\maketitle\tableofcontents
\section{Grundidee}
Nicht jede Akkusativgruppe ist ein direktes Objekt. Bestimmte Akkusativgruppen wirken \textbf{adverbial}: Sie geben Dauer, Zeitpunkt, Ausdehnung oder Maß an, ohne vom Verb als Akkusativobjekt gefordert zu sein.
\textbf{English:} Some accusative phrases function as adverbials (duration, extent, measure) rather than direct objects.

\section{Akkusativ der Zeit (temporaler Akkusativ)}
Er beantwortet insbesondere \textit{Wie lange?} oder \textit{Wann?}.
\begin{itemize}
\item \textit{Ich habe den ganzen Tag gelernt.} -- I studied all day. (Dauer)
\item \textit{Wir bleiben einen Monat in Graz.} -- We are staying in Graz for a month. (Dauer)
\item \textit{Nächsten Freitag habe ich frei.} -- I am off next Friday. (Zeitpunkt)
\item \textit{Jeden Morgen trinkt sie Tee.} -- She drinks tea every morning. (Wiederholung)
\end{itemize}
Auch ohne Präposition sind diese Zeitangaben möglich. Vergleich: \textit{seit einem Monat} enthält dagegen \textit{seit + Dativ} und ist kein freier Akkusativ.

\section{Akkusativ des Maßes und der Ausdehnung}
Er gibt an, \textbf{wie groß, lang, weit, schwer} eine Menge oder ein Unterschied ist.
\begin{itemize}
\item \textit{Der Tisch ist einen Meter lang.} -- The table is one metre long.
\item \textit{Sie ist einen Kopf größer als ich.} -- She is a head taller than I am.
\item \textit{Wir sind einen Kilometer gelaufen.} -- We walked one kilometre.
\item \textit{Das Paket wiegt zwei Kilogramm.} -- The parcel weighs two kilograms.
\end{itemize}
Bei Maßangaben hängt die genaue syntaktische Analyse teilweise vom Verb und der Konstruktion ab. Entscheidend für Lernende ist die Unterscheidung zur gewöhnlichen Gegenstands-Ergänzung.

\section{Unterschied zum Akkusativobjekt}
\textit{Ich lese das Buch.} -- \textbf{Was} lese ich? \textit{das Buch} = Akkusativobjekt.\\
\textit{Ich lese den ganzen Abend.} -- \textbf{Wie lange} lese ich? \textit{den ganzen Abend} = adverbialer Akkusativ.

\textit{Ich kaufe einen Meter Stoff.} -- Die ganze Nominalgruppe bezeichnet die gekaufte Menge; das ist nicht dieselbe Funktion wie in \textit{Der Stoff ist einen Meter breit.}

\section{Keine direkte Entsprechung zum Dativus commodi}
Der \textbf{Dativus commodi} nennt eine \emph{begünstigte Person}: \textit{Ich kaufe mir eine Jacke.}\\
Der \textbf{freie/adverbiale Akkusativ} nennt typischerweise \emph{Zeit oder Maß}: \textit{Ich arbeite den ganzen Tag.}\\
Es gibt keine gleichartige, allgemein produktive Akkusativkonstruktion, die einfach den Begünstigten bezeichnet. \textit{für mich} ist eine Präpositionalgruppe mit \textit{für + Akkusativ}, \emph{kein freier Akkusativ}.

\section{Akkusativus relationis: Terminologie}
Der traditionelle Ausdruck \textit{Akkusativus relationis} bezeichnet einen Akkusativ hinsichtlich eines Bezugsbereichs, vor allem in historischen oder vergleichenden Grammatiktraditionen. Er ist \textbf{nicht} einfach der Standardname aller freien Akkusative im heutigen Deutsch. Beispielsweise wird \textit{einen Kopf} in \textit{einen Kopf größer} in der Gegenwartsgrammatik sinnvoll als \textbf{Differenz- bzw.\ Maßakkusativ} erklärt.

\section{Kasus erkennen}
\begin{itemize}
\item \textit{einen Tag} (Akk.) gegenüber \textit{einem Tag} (Dat.).
\item \textit{den ganzen Tag} (Akk.) gegenüber \textit{an dem ganzen Tag} (Präpositionalgruppe; Dat.).
\item \textit{einen Meter} (Akk.) gegenüber \textit{mit einem Meter} (\textit{mit + Dat.}).
\end{itemize}

\section{Übung}
Unterscheide Akkusativobjekt und adverbialen Akkusativ:\\
(1) \textit{Ich sehe den Film.} (2) \textit{Ich warte den ganzen Abend.} (3) \textit{Der Schrank ist einen Meter breit.} (4) \textit{Ich kaufe den Schrank.}

\textbf{Lösungen:} (1) Objekt; (2) Zeitakkusativ; (3) Maßakkusativ; (4) Objekt.
\end{document}
